% Included by Exercises.tex and Week_4_Exercise_Sheet.tex.
\section{Week 4 --- Theoretical Applications of the It\^o Formula}

Unless otherwise specified, \(W\) is standard Brownian motion, all initial
values are deterministic, and coefficients are constant. Work on bounded
time horizons and justify each use of a zero-mean stochastic integral.

\subsection*{Part I --- Transformations and Their Meaning}

\begin{exercise}[Easy --- A transform whose curvature changes sign]
Let \(Y_t=\log(1+W_t^2)\).
\begin{exparts}
  \item For \(f(x)=\log(1+x^2)\), compute \(f'(x)\) and \(f''(x)\).
  Explain why the classical It\^o formula applies on the whole real line.
  \item Find the drift and diffusion coefficient of \(Y\).
  Determine the regions where its drift is positive, zero, or negative.
  \item Find an explicit function \(a\) such that
  \[
  Y_t-\int_0^t a(W_s)\dd s
  \]
  is a true martingale. Check this by bounding the diffusion coefficient.
  \item Compute \([Y]_t\) as an integral. Explain why the sign changes
  in the drift do not produce negative increments of quadratic variation.
\end{exparts}
\end{exercise}

\begin{exercise}[Easy to medium --- Recovering the noise by an inverse transform]
Set \(X_t=\sinh(W_t)\), and recall that
\[
F(x)=\operatorname{arsinh}(x)=\log\bigl(x+\sqrt{1+x^2}\bigr).
\]
\begin{exparts}
  \item Apply It\^o to \(\sinh(W_t)\) and express both coefficients
  entirely in terms of \(X_t\). Use \(\cosh^2 w-\sinh^2 w=1\).
  \item Compute \(F'(x)\) and \(F''(x)\), then apply It\^o to
  \(F(X_t)\). Check explicitly that the drift cancels and the diffusion
  coefficient is one.
  \item A student keeps only \(F'(X_t)\dd X_t\). Compute the spurious
  drift produced by that calculation.
  \item Now let \(\dd Z_t=a(Z_t)\dd t+\sqrt{1+Z_t^2}\dd W_t\).
  Determine the drift function \(a(x)\) for which the drift of
  \(F(Z_t)\) vanishes identically. In that case, identify \(F(Z_t)\)
  from its initial value and differential.
\end{exparts}
\end{exercise}

\clearpage
\subsection*{Part II --- Which Martingale Claims Survive?}

\begin{exercise}[Easy to medium --- The same transform under two dynamics]
Let \(f\in C^2(\R)\) and fix \(\mu\neq0\). Consider, for every starting
point \(x\in\R\),
\[
X_t^x=x+W_t,
\qquad
Y_t^x=x+\mu t+W_t.
\]
\begin{exparts}
  \item Write the drift coefficient of \(f(X_t^x)\) and of
  \(f(Y_t^x)\).
  \item Determine all \(f\) for which the drift of \(f(X_t^x)\)
  vanishes identically for every \(x\).
  \item Among those functions, determine which also give identically
  zero drift for \(f(Y_t^x)\). Explain why a martingale conclusion
  depends on the dynamics of the input process.
  \item Verify directly that the candidate processes in parts (b) and
  (c) are true martingales on bounded horizons, rather than stopping at
  a formal zero-drift calculation.
\end{exparts}
\end{exercise}

\begin{exercise}[Medium --- Correlation between two compensated squares]
Let \((W,B)\) be jointly Brownian with
\(\dd[W,B]_t=\rho\dd t\), \(-1\leq\rho\leq1\), and set
\[
M_t=W_t^2-t,
\qquad
N_t=B_t^2-t.
\]
You may use \(\E[W_tB_t]=\rho t\) and the fact that Gaussian variables
have finite moments of every order.
\begin{exparts}
  \item Find \(\dd M_t\), \(\dd N_t\), and \(\dd[M,N]_t\).
  \item Apply the product rule to \(M_tN_t\). Find the finite-variation
  term that must be subtracted to obtain a true martingale. Justify the
  integrability of its stochastic terms using Gaussian moments.
  \item Take expectations to compute \(\E[M_tN_t]\), then find the
  correlation of \(M_t\) and \(N_t\) for \(t>0\).
  \emph{Hint:} each compensated square has variance \(2t^2\).
  \item Explain why changing \(\rho\) to \(-\rho\) leaves this
  correlation unchanged. For which \(\rho\) is \(MN\) itself a
  martingale? Check the cases \(\rho=0\) and \(|\rho|=1\).
\end{exparts}
\end{exercise}

\clearpage
\subsection*{Part III --- Products with Accumulated Quantities}

\begin{exercise}[Medium --- A random discount factor has finite variation]
Define
\[
A_t=\int_0^t W_s^2\dd s,
\qquad D_t=e^{-A_t},
\qquad Z_t=D_tW_t.
\]
\begin{exparts}
  \item Find \(\dd A_t\) and \(\dd D_t\). Explain why both processes
  have finite variation on bounded intervals and why \(0<D_t\leq1\).
  \item Compute \([D,W]_t\) and apply the product rule to \(Z_t\).
  Explain why randomness of \(D\) does not create a mixed correction here.
  \item Find \(h_t\) explicitly so that
  \[
  K_t=Z_t+\int_0^t h_s\dd s
  \]
  is a true martingale starting at zero. Prove the required expected-energy
  bound for its stochastic integrand.
  \item Compute \([K]_t\) and show that \(\E[K_t^2]\leq t\).
  Does adding back the finite-variation term change quadratic variation?
\end{exparts}
\end{exercise}

\begin{exercise}[Medium --- A geometric model with a continuous inflow]
Suppose an account receives cash at a constant rate \(c>0\) and satisfies
\[
\dd X_t=(\mu X_t+c)\dd t+\sigma X_t\dd W_t,
\qquad X_0=x>0.
\]
Let \(H_t=\exp((\mu-\sigma^2/2)t+\sigma W_t)\), so
\(\dd H_t=\mu H_t\dd t+\sigma H_t\dd W_t\) and \(H_0=1\).
\begin{exparts}
  \item Compute \(\dd(H_t^{-1})\) and \(\dd[X,H^{-1}]_t\).
  Both \(X\) and \(H\) are driven by the same \(W\).
  \item Apply the product rule to \(X_tH_t^{-1}\). Check that its
  Brownian term and all drift terms proportional to \(X_t\) cancel.
  \item Integrate the remaining differential to obtain an explicit
  expression for \(X_t\) in terms of \(H\) and an ordinary time integral.
  Use it to prove \(X_t>0\).
  \item A proposed solution is \(\widehat X_t=H_t(x+ct)\).
  Differentiate it and identify the inflow rate it actually implements.
  Explain why a contribution made at time \(s\) must grow by the factor
  \(H_t/H_s\) between \(s\) and \(t\).
\end{exparts}
\end{exercise}

\subsection*{Part IV --- Intensive It\^o Formula Practice}

\begin{exercise}[Easy to medium --- A time-dependent exponential martingale]
For \(\lambda\in\R\), define
\[
M_t^{\lambda}=\exp\!\left(\lambda W_t-\tfrac12\lambda^2t\right).
\]
\begin{exparts}
  \item Apply the time-dependent It\^o formula to
  \(f(t,x)=\exp(\lambda x-\lambda^2t/2)\), and show that
  \[\dd M_t^{\lambda}=\lambda M_t^{\lambda}\dd W_t.\]
  Identify precisely which two drift contributions cancel.
  \item Prove that \(M^\lambda\) is a true martingale on every bounded
  horizon. You may compute \(\E[(M_t^\lambda)^2]\) from the Gaussian
  moment-generating function.
  \item Let \(a,b\in\R\). Determine all pairs \((a,b)\) for which
  \(\exp(aW_t+bt)\) has zero drift. Explain why the condition is about
  the quadratic variation of \(W\), not ordinary calculus.
  \item Use conditional expectation to prove
  \(\E[M_t^\lambda\mid\F_s]=M_s^\lambda\) for \(s\leq t\), directly
  from independent Brownian increments.
\end{exparts}
\end{exercise}


\begin{exercise}[Medium --- A backward equation is a martingale test]
Let \(u\in C^{1,2}([0,T]\times\R)\) have bounded \(u_x\), and define
\(Z_t=u(t,W_t)\).
\begin{exparts}
  \item Apply It\^o's formula and show that \(Z\) is a true martingale
  whenever
  \[
  u_t+\tfrac12u_{xx}=0.
  \]
  Carefully state where boundedness of \(u_x\) is used.
  \item Let \(g\in C^2(\R)\) have bounded first and second derivatives,
  and set
  \[
  u(t,x)=\E\bigl[g(x+W_{T-t})\bigr].
  \]
  Explain, using independent increments and the martingale property from
  part (a), why \(u(t,W_t)=\E[g(W_T)\mid\F_t]\).
  \item For \(g(x)=\cos x\), compute \(u\) explicitly and verify the
  backward equation by differentiation.
  \item Explain why the same argument cannot apply directly to
  \(g(x)=|x|\) at \(x=0\), and name the extension of It\^o's formula
  that addresses this kink.
\end{exparts}
\end{exercise}
